%% GENERATED by generate_papers.py -- edit the bank_*.py files, not this file.
\documentclass[11pt,a4paper]{article}
\usepackage[a4paper, top=20mm, bottom=20mm, left=18mm, right=18mm, headheight=14pt]{geometry}
\usepackage[T1]{fontenc}
\usepackage{lmodern}
\usepackage{amsmath, amssymb, mathtools}
\usepackage{braket}
\usepackage{array, tabularx, booktabs, multicol}
\usepackage{enumitem}
\usepackage[dvipsnames]{xcolor}
\usepackage[most]{tcolorbox}
\usepackage{needspace}
\usepackage{fancyhdr, lastpage}
\usepackage{pifont}
\usepackage{tikz}
\usepackage{graphicx}
\usepackage{eso-pic}
% SILICO FELLER watermark: official logo, centred, ~8 % opacity, behind all content, on every page
\AddToShipoutPictureBG{\AtPageCenter{\makebox(0,0){\tikz\node[opacity=0.08, inner sep=0pt]{\includegraphics[width=0.72\paperwidth]{../../../assets/branding/silico_feller_logo.png}};}}}
\usepackage[hidelinks]{hyperref}
\newcolumntype{L}{>{\raggedright\arraybackslash}X}
\definecolor{brand}{HTML}{1F4E79}
\definecolor{okgreen}{HTML}{1B7F3B}
\setlength{\parindent}{0pt}
\setlength{\parskip}{3pt}
\newcommand{\optlabel}[1]{\textbf{(#1)}}
\newcommand{\correctlabel}[1]{\textcolor{okgreen}{\textbf{(#1)}\,\ding{51}}}
\newenvironment{qitem}[1]{\begin{list}{}{\setlength{\leftmargin}{3.3em}\setlength{\labelwidth}{3.1em}%
  \setlength{\labelsep}{0.2em}\setlength{\topsep}{3pt}\setlength{\parsep}{3pt}}\item[\textbf{Q.#1}\hfill]}{\end{list}}
\newcommand{\qsep}{\par\vspace{2pt}{\color{black!15}\hrule}\vspace{4pt}}
\newcommand{\qmeta}[2]{{\footnotesize\color{black!60}\textsc{#1} \;$\cdot$\; #2}\par}
\newtcolorbox{passagebox}[1]{enhanced, breakable, colback=brand!4, colframe=brand!70, fonttitle=\bfseries,
  title={#1}, left=6pt, right=6pt, top=4pt, bottom=4pt}
\newtcolorbox{sectionbanner}{enhanced, colback=brand, colframe=brand, coltext=white, fontupper=\large\bfseries,
  left=8pt, right=8pt, top=6pt, bottom=6pt, sharp corners}
\newtcolorbox{infobox}[1]{enhanced, breakable, colback=black!3, colframe=black!45, fonttitle=\bfseries,
  title={#1}, left=6pt, right=6pt, top=4pt, bottom=4pt}
\newtcolorbox{solbox}[1]{enhanced, breakable, colback=okgreen!3, colframe=okgreen!60!black, fonttitle=\bfseries,
  title={Answer: (#1) \quad Detailed solution}, left=6pt, right=6pt, top=3pt, bottom=4pt, before skip=4pt, after skip=6pt}
\newcommand{\HHLop}{\mathrm{HHL}}
\allowdisplaybreaks
\setlength{\emergencystretch}{3em}

\begin{document}
\pagestyle{fancy}\fancyhf{}
\fancyhead[L]{\small\color{black!70}HHL Quantum Linear-Systems Algorithm --- Assessment}
\fancyhead[R]{\small\color{black!70}\textbf{Question Paper}}
\fancyfoot[C]{\small Page \thepage\ of \pageref*{LastPage}}
\renewcommand{\headrulewidth}{0.4pt}

\thispagestyle{empty}
\begin{center}
{\color{brand}\rule{\linewidth}{1.2pt}}\\[6pt]
{\Huge\bfseries\color{brand} HHL Quantum Linear-Systems Algorithm}\\[4pt]
{\Large Assessment of the Harrow--Hassidim--Lloyd algorithm}\\[4pt]
{\large\bfseries QUESTION PAPER}\\[2pt]
{\color{brand}\rule{\linewidth}{1.2pt}}\\[1pt]
\end{center}
\vspace{-2pt}
\begin{tabularx}{\linewidth}{@{}lXlr@{}}
\textbf{Time allowed:} & 3 hours & \textbf{Maximum marks:} & 240\\
\textbf{Questions:} & 80 (4 sections $\times$ 20) & \textbf{Answer key:} & separate Answer Book
\end{tabularx}
\vspace{6pt}

\textbf{Name of the candidate:} \hrulefill\quad \textbf{Roll No.:} \rule{35mm}{0.4pt}

\vspace{6pt}
\begin{infobox}{General instructions}
\begin{enumerate}[leftmargin=1.4em, itemsep=1pt, topsep=1pt]
\item This question paper has \textbf{four sections} of \textbf{20 questions} each:
\textbf{A}~Linear Algebra and Mathematical Foundations, \textbf{B}~Quantum Computing Foundations, \textbf{C}~HHL Algorithm Fundamentals, \textbf{D}~Advanced HHL and Implementation.
\item Every question is of the \textbf{single-correct} type: four options, exactly one correct. Question formats used:
single-correct MCQ, matching lists (List-I/List-II), Assertion--Reason, Statement~I/II and paragraph-based questions.
\item For every question choose \textbf{ONE} option, (A), (B), (C) or (D), and record it against the question number in the answer script provided by the invigilator.
\item Calculators are \textbf{not} allowed. Rough work may be done on separate sheets.
\item Unless a question says otherwise, use the \textbf{conventions} printed below; they apply to every section.
\item The detailed solutions are in the separate \emph{Answer Book}; do not consult it before finishing the test.
\end{enumerate}
\end{infobox}
\begin{infobox}{Marking scheme (every question)}
Each question has \textbf{four} options (A), (B), (C) and (D). \textbf{ONLY ONE} of these four options is correct.\\[2pt]
\begin{tabular}{@{}ll@{}}
\emph{Full marks}: & $+3$ if ONLY the correct option is chosen;\\
\emph{Zero marks}: & $\phantom{+}0$ if none of the options is chosen (question unanswered);\\
\emph{Negative marks}: & $-1$ in all other cases (wrong option, or more than one option chosen).
\end{tabular}
\end{infobox}

\begin{infobox}{Conventions and notation used throughout the paper}
\small
\begin{itemize}[leftmargin=1.2em, itemsep=1pt, topsep=2pt]
\item Linear system $Ax = b$, $A$ a $d\times d$ matrix; $\ket b = b/\|b\|$, $\ket x = x/\|x\|$. For Hermitian $A$: $A = \sum_j\lambda_j\ket{u_j}\bra{u_j}$, $\beta_j = \braket{u_j|b}$;
      $\kappa = |\lambda|_{\max}/|\lambda|_{\min}$ (for general matrices $\sigma_{\max}/\sigma_{\min}$); sparsity $s$ = maximal number of non-zeros per row.
\item Basis states are written with the \textbf{most significant bit first}; $\mathrm{QFT}\ket k = 2^{-n/2}\sum_y e^{2\pi iky/2^n}\ket y$; $R_Y(\theta) = e^{-i\theta Y/2}$.
\item \textbf{Phase estimation} with $t$ counting qubits on $U\ket u = e^{2\pi i\varphi}\ket u$: Hadamards, controlled $U^{2^j}$, inverse QFT; outcome $y \approx 2^t\varphi$.
\item \textbf{HHL in the notebook} (\texttt{HHL.ipynb}): $n_c$ clock qubits, $n_b = \log_2d$ system qubits, 1 ancilla; $U = e^{iAt}$; clock value of $\lambda$:
      $k = \lambda t\,2^{n_c}/(2\pi)$, signed (values $k \ge 2^{n_c-1}$ mean $k - 2^{n_c}$); $C = 2\pi/(2^{n_c}t)$; inversion $R_Y(2\arcsin(C/\lambda_k))$ on the ancilla;
      success = ancilla $1$ and clock $0$, with $P_{\text{success}} = C^2\|A^{-1}b\|^2/\|b\|^2$.
\item \textbf{Default example:} $A = H\operatorname{diag}(1, 2, 3, 4)H$ with $H = \tfrac12\begin{psmallmatrix}1&1&1&1\\1&-1&1&-1\\1&1&-1&-1\\1&-1&-1&1\end{psmallmatrix}$, $b = (1, 0, 0, 0)$, $n_c = 4$, $t = \pi/8$, $C = 1$,
      7 qubits. A \emph{previously recorded run} means the outputs stored in the notebook (CUDA-Q 0.16.0, CPU target \texttt{qpp-cpu}).
\item Numerical answers are rounded.
\end{itemize}
\end{infobox}
\newpage

\begin{sectionbanner}SECTION A \enspace$\cdot$\enspace Linear Algebra and Mathematical Foundations \hfill Q.1--20\end{sectionbanner}

{\small \textbf{Syllabus:} linear systems, eigen-decomposition, Hermitian, unitary and positive-definite matrices, matrix functions, condition numbers, sparsity, normalisation and classical solvers.\\ \textbf{Formats:} Single-correct MCQ (Q.1--19); Matching lists (Q.3--4); Assertion--Reason (Q.7--20); Statement I / II (Q.8); Paragraph based (Q.14--15). Each question: $+3$ / $0$ / $-1$.}
\par\medskip

\needspace{7\baselineskip}
\begin{qitem}{1}
The solution of $\begin{pmatrix}3 & 1\\ 1 & 3\end{pmatrix}x = \begin{pmatrix}1\\ -1\end{pmatrix}$ is
\par\smallskip\noindent\begin{tabularx}{\linewidth}{@{}*{2}{L}@{}}
\optlabel{A}\,$x = (0.5,\ 0.5)^{\mathsf T}$ & \optlabel{B}\,$x = (0.5,\ -0.5)^{\mathsf T}$ \\[3pt]
\optlabel{C}\,$x = (0.25,\ -0.25)^{\mathsf T}$ & \optlabel{D}\,$x = (1,\ -1)^{\mathsf T}$
\end{tabularx}
\end{qitem}
\qsep

\needspace{7\baselineskip}
\begin{qitem}{2}
The eigenvalues and the condition number $\kappa = |\lambda|_{\max}/|\lambda|_{\min}$ of $A = \begin{pmatrix}1 & -\tfrac13\\ -\tfrac13 & 1\end{pmatrix}$ are
\par\smallskip\noindent\begin{tabularx}{\linewidth}{@{}*{2}{L}@{}}
\optlabel{A}\,$\tfrac23$ and $\tfrac43$; $\kappa = \tfrac23$ & \optlabel{B}\,$\tfrac89$ and $1$; $\kappa = \tfrac98$ \\[3pt]
\optlabel{C}\,$1$ and $-\tfrac13$; $\kappa = 3$ & \optlabel{D}\,$\tfrac23$ and $\tfrac43$; $\kappa = 2$
\end{tabularx}
\end{qitem}
\qsep

\needspace{7\baselineskip}
\begin{qitem}{3}
Match each class of square matrix (List-I) with the property of its eigenvalues (List-II).
\par\smallskip\noindent{\renewcommand{\arraystretch}{1.35}%
\begin{tabularx}{\linewidth}{@{}c>{\raggedright\arraybackslash}p{0.33\linewidth}@{\hspace{1.2em}}cL@{}}
\toprule
\multicolumn{2}{@{}l}{\textbf{List-I}} & \multicolumn{2}{l@{}}{\textbf{List-II}} \\ \midrule
\textbf{(P)} & Hermitian ($A = A^\dagger$) & \textbf{(1)} & at least one eigenvalue is $0$ \\
\textbf{(Q)} & unitary ($U^\dagger U = I$) & \textbf{(2)} & all eigenvalues are real \\
\textbf{(R)} & Hermitian positive definite & \textbf{(3)} & all eigenvalues lie on the unit circle, $|\lambda| = 1$ \\
\textbf{(S)} & singular & \textbf{(4)} & all eigenvalues are real and strictly positive \\
\bottomrule
\end{tabularx}}
\par\smallskip\noindent\begin{tabularx}{\linewidth}{@{}*{2}{L}@{}}
\optlabel{A}\,P$\to$2, Q$\to$4, R$\to$3, S$\to$1 & \optlabel{B}\,P$\to$2, Q$\to$3, R$\to$1, S$\to$4 \\[3pt]
\optlabel{C}\,P$\to$2, Q$\to$3, R$\to$4, S$\to$1 & \optlabel{D}\,P$\to$3, Q$\to$2, R$\to$4, S$\to$1
\end{tabularx}
\end{qitem}
\qsep

\needspace{7\baselineskip}
\begin{qitem}{4}
Match each matrix (List-I) with its condition number $\kappa = |\lambda|_{\max}/|\lambda|_{\min}$ (List-II).
\par\smallskip\noindent{\renewcommand{\arraystretch}{1.35}%
\begin{tabularx}{\linewidth}{@{}c>{\raggedright\arraybackslash}p{0.33\linewidth}@{\hspace{1.2em}}cL@{}}
\toprule
\multicolumn{2}{@{}l}{\textbf{List-I}} & \multicolumn{2}{l@{}}{\textbf{List-II}} \\ \midrule
\textbf{(P)} & $\operatorname{diag}(0.5,\ 2,\ 8)$ & \textbf{(1)} & $4$ \\
\textbf{(Q)} & $\begin{psmallmatrix}1 & -1/3\\ -1/3 & 1\end{psmallmatrix}$ & \textbf{(2)} & $1$ \\
\textbf{(R)} & the notebook's default $A$ (eigenvalues $1, 2, 3, 4$) & \textbf{(3)} & $16$ \\
\textbf{(S)} & $\begin{psmallmatrix}2 & 0\\ 0 & 2\end{psmallmatrix}$ & \textbf{(4)} & $2$ \\
\bottomrule
\end{tabularx}}
\par\smallskip\noindent\begin{tabularx}{\linewidth}{@{}*{2}{L}@{}}
\optlabel{A}\,P$\to$1, Q$\to$4, R$\to$3, S$\to$2 & \optlabel{B}\,P$\to$3, Q$\to$4, R$\to$2, S$\to$1 \\[3pt]
\optlabel{C}\,P$\to$3, Q$\to$2, R$\to$1, S$\to$4 & \optlabel{D}\,P$\to$3, Q$\to$4, R$\to$1, S$\to$2
\end{tabularx}
\end{qitem}
\qsep

\needspace{7\baselineskip}
\begin{qitem}{5}
The notebook's default matrix is $A = H\,\operatorname{diag}(1, 2, 3, 4)\,H$ with the symmetric orthogonal matrix
$H = \tfrac12\begin{psmallmatrix}1&1&1&1\\1&-1&1&-1\\1&1&-1&-1\\1&-1&-1&1\end{psmallmatrix}$ ($H^2 = I$). For $b = (1, 0, 0, 0)^{\mathsf T}$, the first component $x_0$ of $x = A^{-1}b$ is
\par\smallskip\noindent\begin{tabularx}{\linewidth}{@{}*{4}{L}@{}}
\optlabel{A}\,$\tfrac{5}{8}$ & \optlabel{B}\,$\tfrac14$ & \optlabel{C}\,$\tfrac{25}{12} \approx 2.0833$ & \optlabel{D}\,$\tfrac{25}{48} \approx 0.5208$
\end{tabularx}
\end{qitem}
\qsep

\needspace{7\baselineskip}
\begin{qitem}{6}
HHL needs a Hermitian matrix, so the notebook embeds a non-Hermitian $A$ as $\tilde A = \begin{pmatrix}0 & A\\ A^\dagger & 0\end{pmatrix}$. For $A = \begin{pmatrix}1 & 2\\ 0 & 1\end{pmatrix}$ the
condition number of $\tilde A$ is
\par\smallskip\noindent\begin{tabularx}{\linewidth}{@{}*{2}{L}@{}}
\optlabel{A}\,$3 + 2\sqrt2 \approx 5.83$ & \optlabel{B}\,$\sqrt2 + 1 \approx 2.41$ \\[3pt]
\optlabel{C}\,$1$, because both eigenvalues of $A$ are $1$ & \optlabel{D}\,$2$
\end{tabularx}
\end{qitem}
\qsep

\needspace{7\baselineskip}
\begin{qitem}{7}
\textbf{Assertion (A):} Every invertible matrix has condition number $\kappa(A) = \|A\|\,\|A^{-1}\| \ge 1$.\\[1mm]
\textbf{Reason (R):} In the spectral norm $\kappa(A) = \sigma_{\max}/\sigma_{\min}$, the ratio of the largest to the smallest singular value.
\begin{itemize}[leftmargin=2.6em, labelsep=0.4em, itemsep=2pt, topsep=3pt]
\item[\optlabel{A}] Both Assertion (A) and Reason (R) are true, and R is the correct explanation of A
\item[\optlabel{B}] Both A and R are true, but R is \emph{not} the correct explanation of A
\item[\optlabel{C}] A is true, but R is false
\item[\optlabel{D}] A is false, but R is true
\end{itemize}
\end{qitem}
\qsep

\needspace{7\baselineskip}
\begin{qitem}{8}
\textbf{Statement I:} Every invertible matrix is Hermitian.\\[1mm]
\textbf{Statement II:} Every real symmetric matrix can be written as $A = Q\Lambda Q^{\mathsf T}$ with an orthogonal $Q$ and a real diagonal $\Lambda$.
\begin{itemize}[leftmargin=2.6em, labelsep=0.4em, itemsep=2pt, topsep=3pt]
\item[\optlabel{A}] Both Statement I and Statement II are true
\item[\optlabel{B}] Both Statement I and Statement II are false
\item[\optlabel{C}] Statement I is true, but Statement II is false
\item[\optlabel{D}] Statement I is false, but Statement II is true
\end{itemize}
\end{qitem}
\qsep

\needspace{7\baselineskip}
\begin{qitem}{9}
The sparsity $s$ of a matrix is the largest number of non-zero entries in any row. For the notebook's default $4\times4$ matrix
$A = \begin{psmallmatrix}2.5&-0.5&-1&0\\-0.5&2.5&0&-1\\-1&0&2.5&-0.5\\0&-1&-0.5&2.5\end{psmallmatrix}$ and for a tridiagonal $1000\times1000$ matrix, $s$ is
\par\smallskip\noindent\begin{tabularx}{\linewidth}{@{}*{4}{L}@{}}
\optlabel{A}\,$3$ and $1000$ & \optlabel{B}\,$12$ and $2998$ & \optlabel{C}\,$3$ and $3$ & \optlabel{D}\,$4$ and $3$
\end{tabularx}
\end{qitem}
\qsep

\needspace{7\baselineskip}
\begin{qitem}{10}
Gaussian elimination on a dense $d\times d$ system costs about $\tfrac23d^3$ floating-point operations. Doubling $d$ from $1000$ to $2000$ multiplies the work by
\par\smallskip\noindent\begin{tabularx}{\linewidth}{@{}*{4}{L}@{}}
\optlabel{A}\,$8$ & \optlabel{B}\,$2$ & \optlabel{C}\,$4$ & \optlabel{D}\,$\tfrac23$
\end{tabularx}
\end{qitem}
\qsep

\needspace{7\baselineskip}
\begin{qitem}{11}
For a sparse symmetric positive-definite matrix with condition number $\kappa$, the number of conjugate-gradient iterations needed for a fixed accuracy
scales as
\par\smallskip\noindent\begin{tabularx}{\linewidth}{@{}*{4}{L}@{}}
\optlabel{A}\,$O(\log\kappa)$ & \optlabel{B}\,$O(d^3)$ & \optlabel{C}\,$O(\sqrt\kappa)$ & \optlabel{D}\,$O(\kappa^2)$
\end{tabularx}
\end{qitem}
\qsep

\needspace{7\baselineskip}
\begin{qitem}{12}
The vector $b = (1, 2, 2, 0)^{\mathsf T}$ is amplitude-encoded on two qubits as $\ket b = \sum_ib_i\ket i/\|b\|$ (index $i$ in binary, first qubit most significant). The
probability of measuring $\ket{01}$ is
\par\smallskip\noindent\begin{tabularx}{\linewidth}{@{}*{4}{L}@{}}
\optlabel{A}\,$\tfrac13$ & \optlabel{B}\,$\tfrac29$ & \optlabel{C}\,$\tfrac23$ & \optlabel{D}\,$\tfrac49$
\end{tabularx}
\end{qitem}
\qsep

\needspace{7\baselineskip}
\begin{qitem}{13}
For the notebook's default $A = H\operatorname{diag}(1, 2, 3, 4)H$ (eigenvectors = columns of $H$) and $b = (1, 0, 0, 0)$, the coefficients $\beta_j = \braket{u_j|b}$ of
$\ket b$ in the eigenbasis are
\par\smallskip\noindent\begin{tabularx}{\linewidth}{@{}*{2}{L}@{}}
\optlabel{A}\,$(1, 0, 0, 0)$ & \optlabel{B}\,$(\tfrac12, \tfrac12, \tfrac12, \tfrac12)$ \\[3pt]
\optlabel{C}\,$(\tfrac12, -\tfrac12, \tfrac12, -\tfrac12)$ & \optlabel{D}\,$(\tfrac14, \tfrac14, \tfrac14, \tfrac14)$
\end{tabularx}
\end{qitem}
\qsep

\needspace{20\baselineskip}
\begin{passagebox}{Paragraph for Questions 14 and 15}
The notebook's 2$\times$2 example is $A = \begin{pmatrix}1 & -\tfrac13\\ -\tfrac13 & 1\end{pmatrix}$, $b = \begin{pmatrix}1\\0\end{pmatrix}$. Its eigenvectors are
$u_1 = \tfrac{1}{\sqrt2}(1, 1)^{\mathsf T}$ and $u_2 = \tfrac{1}{\sqrt2}(1, -1)^{\mathsf T}$.
\end{passagebox}
\needspace{7\baselineskip}
\begin{qitem}{14}
The solution $x = A^{-1}b$ is
\par\smallskip\noindent\begin{tabularx}{\linewidth}{@{}*{4}{L}@{}}
\optlabel{A}\,$(\tfrac98,\ \tfrac38)^{\mathsf T}$ & \optlabel{B}\,$(1,\ \tfrac13)^{\mathsf T}$ & \optlabel{C}\,$(\tfrac34,\ \tfrac14)^{\mathsf T}$ & \optlabel{D}\,$(\tfrac98,\ -\tfrac38)^{\mathsf T}$
\end{tabularx}
\end{qitem}
\qsep

\needspace{7\baselineskip}
\begin{qitem}{15}
HHL prepares the normalised state $\ket x = x/\|x\|$. Its amplitudes, and the probability of measuring $\ket0$ on the system qubit, are
\par\smallskip\noindent\begin{tabularx}{\linewidth}{@{}*{2}{L}@{}}
\optlabel{A}\,$\tfrac12(\sqrt3, 1)$ and $0.75$ & \optlabel{B}\,$(\tfrac98, \tfrac38)$ and $\tfrac98$ \\[3pt]
\optlabel{C}\,$\tfrac{1}{\sqrt{10}}(3, 1)$ and $0.9$ & \optlabel{D}\,$\tfrac{1}{\sqrt{10}}(3, 1)$ and $0.75$
\end{tabularx}
\end{qitem}
\qsep

\needspace{7\baselineskip}
\begin{qitem}{16}
For $A = \operatorname{diag}(1, 2)$ and $t = \pi/2$, the matrix $e^{iAt}$ is
\par\smallskip\noindent\begin{tabularx}{\linewidth}{@{}*{4}{L}@{}}
\optlabel{A}\,$\operatorname{diag}(i,\ 1)$ & \optlabel{B}\,$\operatorname{diag}(i,\ -1)$ & \optlabel{C}\,$\begin{pmatrix}0 & i\\ i & 0\end{pmatrix}$ & \optlabel{D}\,$\operatorname{diag}(-i,\ -1)$
\end{tabularx}
\end{qitem}
\qsep

\needspace{7\baselineskip}
\begin{qitem}{17}
For which matrices $A$ is $e^{iAt}$ unitary for every real $t$?
\par\smallskip\noindent\begin{tabularx}{\linewidth}{@{}*{2}{L}@{}}
\optlabel{A}\,exactly the sparse matrices & \optlabel{B}\,exactly the Hermitian matrices \\[3pt]
\optlabel{C}\,exactly the invertible matrices & \optlabel{D}\,exactly the real matrices
\end{tabularx}
\end{qitem}
\qsep

\needspace{7\baselineskip}
\begin{qitem}{18}
For the default system ($A$ with eigenvalues $1, 2, 3, 4$ and $\|b\| = 1$), the norm $\|x\| = \|A^{-1}b\|$ is $0.5966$. Which bound must every such norm satisfy?
\begin{itemize}[leftmargin=2.6em, labelsep=0.4em, itemsep=2pt, topsep=3pt]
\item[\optlabel{A}] $\dfrac{\|b\|}{\lambda_{\max}} \le \|x\| \le \dfrac{\|b\|}{\lambda_{\min}}$, i.e.\ $0.25 \le 0.5966 \le 1$
\item[\optlabel{B}] $\|x\| = \|b\|/\kappa = 0.25$ exactly
\item[\optlabel{C}] $\|x\| \le \|b\|/\lambda_{\max} = 0.25$
\item[\optlabel{D}] $\|b\|\,\lambda_{\min} \le \|x\| \le \|b\|\,\lambda_{\max}$, i.e.\ $1 \le \|x\| \le 4$
\end{itemize}
\end{qitem}
\qsep

\needspace{7\baselineskip}
\begin{qitem}{19}
For which values of $a$ is $\begin{pmatrix}a & 1\\ 1 & a\end{pmatrix}$ singular (so that HHL cannot be applied)?
\par\smallskip\noindent\begin{tabularx}{\linewidth}{@{}*{4}{L}@{}}
\optlabel{A}\,$a = \pm2$ & \optlabel{B}\,$a = \pm1$ & \optlabel{C}\,$a = 0$ & \optlabel{D}\,no real $a$
\end{tabularx}
\end{qitem}
\qsep

\needspace{7\baselineskip}
\begin{qitem}{20}
\textbf{Assertion (A):} Multiplying an invertible matrix by a non-zero constant $c$ does not change its condition number.\\[1mm]
\textbf{Reason (R):} Multiplying $A$ by $c$ multiplies every eigenvalue by $c^2$.
\begin{itemize}[leftmargin=2.6em, labelsep=0.4em, itemsep=2pt, topsep=3pt]
\item[\optlabel{A}] Both Assertion (A) and Reason (R) are true, and R is the correct explanation of A
\item[\optlabel{B}] Both A and R are true, but R is \emph{not} the correct explanation of A
\item[\optlabel{C}] A is true, but R is false
\item[\optlabel{D}] A is false, but R is true
\end{itemize}
\end{qitem}
\qsep

\newpage

\begin{sectionbanner}SECTION B \enspace$\cdot$\enspace Quantum Computing Foundations \hfill Q.21--40\end{sectionbanner}

{\small \textbf{Syllabus:} amplitude encoding, measurement and post-selection, controlled gates and phase kickback, QFT and inverse QFT, quantum phase estimation, two's complement clock values, rotations, entanglement and memory.\\ \textbf{Formats:} Single-correct MCQ (Q.21--40); Matching lists (Q.23--30); Assertion--Reason (Q.35); Statement I / II (Q.29--39); Paragraph based (Q.37--38). Each question: $+3$ / $0$ / $-1$.}
\par\medskip

\needspace{7\baselineskip}
\begin{qitem}{21}
Phase estimation with 4 counting qubits is applied to an eigenstate of $U$ with eigenvalue $e^{2\pi i\cdot3/16}$. The measured counting register is
\begin{itemize}[leftmargin=2.6em, labelsep=0.4em, itemsep=2pt, topsep=3pt]
\item[\optlabel{A}] $\ket{0011}$ with probability $\tfrac{3}{16}$
\item[\optlabel{B}] $\ket{0011}$ or $\ket{0100}$, each with probability $\tfrac12$
\item[\optlabel{C}] $\ket{0011}$ with probability 1
\item[\optlabel{D}] $\ket{1100}$ with probability 1
\end{itemize}
\end{qitem}
\qsep

\needspace{7\baselineskip}
\begin{qitem}{22}
In HHL, phase estimation uses $U = e^{iAt}$, so an eigenvalue $\lambda$ gives the clock value $k = \lambda t\,2^{n_c}/(2\pi)$. With $n_c = 4$ and $t = \pi/8$ (the notebook's
default), the eigenvalue $\lambda = 4$ gives
\par\smallskip\noindent\begin{tabularx}{\linewidth}{@{}*{2}{L}@{}}
\optlabel{A}\,$k = 2$ & \optlabel{B}\,$k = 4\pi$ \\[3pt]
\optlabel{C}\,$k = 8$ & \optlabel{D}\,$k = 4$, i.e.\ $\ket{0100}$
\end{tabularx}
\end{qitem}
\qsep

\needspace{7\baselineskip}
\begin{qitem}{23}
Match each step of quantum phase estimation (List-I) with its effect (List-II).
\par\smallskip\noindent{\renewcommand{\arraystretch}{1.35}%
\begin{tabularx}{\linewidth}{@{}c>{\raggedright\arraybackslash}p{0.33\linewidth}@{\hspace{1.2em}}cL@{}}
\toprule
\multicolumn{2}{@{}l}{\textbf{List-I}} & \multicolumn{2}{l@{}}{\textbf{List-II}} \\ \midrule
\textbf{(P)} & Hadamards on the counting register & \textbf{(1)} & converts the pattern of phases $e^{2\pi i2^j\varphi}$ into the binary number $\approx 2^t\varphi$ \\
\textbf{(Q)} & controlled $U^{2^j}$ gates & \textbf{(2)} & returns one estimate $y$ of $2^t\varphi$ \\
\textbf{(R)} & inverse QFT on the counting register & \textbf{(3)} & creates the uniform superposition $\frac{1}{\sqrt{2^t}}\sum_j\ket j$ \\
\textbf{(S)} & measurement of the counting register & \textbf{(4)} & write the phase $e^{2\pi i2^j\varphi}$ onto counting qubit $j$ (phase kickback) \\
\bottomrule
\end{tabularx}}
\par\smallskip\noindent\begin{tabularx}{\linewidth}{@{}*{2}{L}@{}}
\optlabel{A}\,P$\to$3, Q$\to$1, R$\to$4, S$\to$2 & \optlabel{B}\,P$\to$3, Q$\to$4, R$\to$2, S$\to$1 \\[3pt]
\optlabel{C}\,P$\to$4, Q$\to$3, R$\to$1, S$\to$2 & \optlabel{D}\,P$\to$3, Q$\to$4, R$\to$1, S$\to$2
\end{tabularx}
\end{qitem}
\qsep

\needspace{7\baselineskip}
\begin{qitem}{24}
A control qubit in $\ket+$ and a target in an eigenstate $\ket u$ of $U$ ($U\ket u = e^{i\phi}\ket u$) pass through a controlled-$U$. The result is
\begin{itemize}[leftmargin=2.6em, labelsep=0.4em, itemsep=2pt, topsep=3pt]
\item[\optlabel{A}] $\tfrac{1}{\sqrt2}(\ket0 + e^{i\phi}\ket1)\otimes\ket u$
\item[\optlabel{B}] $\tfrac{1}{\sqrt2}(\ket0\ket u + \ket1\ket{u'})$ with $\ket{u'}$ orthogonal to $\ket u$
\item[\optlabel{C}] $\tfrac{1}{\sqrt2}(e^{i\phi}\ket0 + \ket1)\otimes\ket u$
\item[\optlabel{D}] $\ket+\otimes e^{i\phi}\ket u$
\end{itemize}
\end{qitem}
\qsep

\needspace{7\baselineskip}
\begin{qitem}{25}
The inverse QFT on two qubits is applied to $\tfrac12\big(\ket{00} + i\ket{01} - \ket{10} - i\ket{11}\big)$. The result is
\par\smallskip\noindent\begin{tabularx}{\linewidth}{@{}*{2}{L}@{}}
\optlabel{A}\,$\tfrac12(\ket{00} + \ket{01} + \ket{10} + \ket{11})$ & \optlabel{B}\,$\ket{01}$ \\[3pt]
\optlabel{C}\,$\ket{11}$ & \optlabel{D}\,$\ket{10}$
\end{tabularx}
\end{qitem}
\qsep

\needspace{7\baselineskip}
\begin{qitem}{26}
The notebook builds the inverse QFT (in phase estimation) and the QFT (in the uncomputation) on $n_c = 4$ clock qubits from controlled phase gates
\texttt{cr1} and \texttt{swap}s. How many \texttt{cr1} and \texttt{swap} gates does the whole HHL circuit contain?
\par\smallskip\noindent\begin{tabularx}{\linewidth}{@{}*{4}{L}@{}}
\optlabel{A}\,$12$ and $8$ & \optlabel{B}\,$16$ and $4$ & \optlabel{C}\,$6$ and $2$ & \optlabel{D}\,$12$ and $4$
\end{tabularx}
\end{qitem}
\qsep

\needspace{7\baselineskip}
\begin{qitem}{27}
An ancilla in $\ket0$ is rotated by $R_Y(\theta) = e^{-i\theta Y/2}$ with $\theta = 2\arcsin(\tfrac12)$. The probability of then measuring $\ket1$ is
\par\smallskip\noindent\begin{tabularx}{\linewidth}{@{}*{4}{L}@{}}
\optlabel{A}\,$\tfrac34$ & \optlabel{B}\,$\tfrac14$ & \optlabel{C}\,$\sin^2(\tfrac12) \approx 0.23$ & \optlabel{D}\,$\tfrac12$
\end{tabularx}
\end{qitem}
\qsep

\needspace{7\baselineskip}
\begin{qitem}{28}
To apply a rotation only when the 4-qubit clock register is in $\ket k$, the notebook surrounds a fully controlled \texttt{cry} with $X$ gates on the clock qubits whose bit of $k$ is
0 (clock[0] = most significant bit). For $k = 5$, the $X$ gates act on
\par\smallskip\noindent\begin{tabularx}{\linewidth}{@{}*{2}{L}@{}}
\optlabel{A}\,clock[2] and clock[3] & \optlabel{B}\,clock[0] and clock[2] \\[3pt]
\optlabel{C}\,clock[0] only & \optlabel{D}\,clock[1] and clock[3]
\end{tabularx}
\end{qitem}
\qsep

\needspace{7\baselineskip}
\begin{qitem}{29}
\textbf{Statement I:} A controlled-$U^{2^j}$ can be implemented as $2^j$ consecutive controlled-$U$ gates.\\[1mm]
\textbf{Statement II:} For $U = e^{iAt}$, $U^{2^j} = e^{iA\,t\,2^j}$.
\begin{itemize}[leftmargin=2.6em, labelsep=0.4em, itemsep=2pt, topsep=3pt]
\item[\optlabel{A}] Both Statement I and Statement II are true
\item[\optlabel{B}] Both Statement I and Statement II are false
\item[\optlabel{C}] Statement I is true, but Statement II is false
\item[\optlabel{D}] Statement I is false, but Statement II is true
\end{itemize}
\end{qitem}
\qsep

\needspace{7\baselineskip}
\begin{qitem}{30}
HHL reads the 4-qubit clock as a signed (two's complement) number: values $k \ge 8$ mean $k - 16$. Match each clock value (List-I) with the signed value
(List-II).
\par\smallskip\noindent{\renewcommand{\arraystretch}{1.35}%
\begin{tabularx}{\linewidth}{@{}c>{\raggedright\arraybackslash}p{0.33\linewidth}@{\hspace{1.2em}}cL@{}}
\toprule
\multicolumn{2}{@{}l}{\textbf{List-I}} & \multicolumn{2}{l@{}}{\textbf{List-II}} \\ \midrule
\textbf{(P)} & $k = 3$ & \textbf{(1)} & $-1$ \\
\textbf{(Q)} & $k = 8$ & \textbf{(2)} & $-4$ \\
\textbf{(R)} & $k = 12$ & \textbf{(3)} & $3$ \\
\textbf{(S)} & $k = 15$ & \textbf{(4)} & $-8$ \\
\bottomrule
\end{tabularx}}
\par\smallskip\noindent\begin{tabularx}{\linewidth}{@{}*{2}{L}@{}}
\optlabel{A}\,P$\to$3, Q$\to$2, R$\to$4, S$\to$1 & \optlabel{B}\,P$\to$3, Q$\to$4, R$\to$1, S$\to$2 \\[3pt]
\optlabel{C}\,P$\to$3, Q$\to$4, R$\to$2, S$\to$1 & \optlabel{D}\,P$\to$1, Q$\to$4, R$\to$2, S$\to$3
\end{tabularx}
\end{qitem}
\qsep

\needspace{7\baselineskip}
\begin{qitem}{31}
A state $\ket\Psi = \sqrt{1 - p}\,\ket{\psi_0}\ket0 + \sqrt p\,\ket{\psi_1}\ket1$ (normalised $\ket{\psi_0}$, $\ket{\psi_1}$) is measured on the last qubit and the result 1 is kept (post-selection).
Which statement is correct?
\begin{itemize}[leftmargin=2.6em, labelsep=0.4em, itemsep=2pt, topsep=3pt]
\item[\optlabel{A}] The outcome 1 occurs with probability $p$ and leaves the rest in $\ket{\psi_1}$
\item[\optlabel{B}] The outcome 1 occurs with probability $\sqrt p$
\item[\optlabel{C}] Post-selection succeeds with probability 1 if $p > 0$
\item[\optlabel{D}] The remaining state is $\sqrt p\,\ket{\psi_1}$, not normalised
\end{itemize}
\end{qitem}
\qsep

\needspace{7\baselineskip}
\begin{qitem}{32}
After phase estimation in HHL the system and clock registers are in $\sum_j\beta_j\ket{u_j}\ket{k_j}$ (distinct $k_j$). When is this a product state of the two registers?
\begin{itemize}[leftmargin=2.6em, labelsep=0.4em, itemsep=2pt, topsep=3pt]
\item[\optlabel{A}] always
\item[\optlabel{B}] never
\item[\optlabel{C}] when $\ket b$ is an eigenvector of $A$ (only one $\beta_j \ne 0$)
\item[\optlabel{D}] when all $\beta_j$ are equal
\end{itemize}
\end{qitem}
\qsep

\needspace{7\baselineskip}
\begin{qitem}{33}
A state vector of $n$ qubits needs $16\times2^n$ bytes (fp64). The notebook's default HHL circuit has 7 qubits, yet it prints ``Estimated memory: 4.10 KB''. Why?
\begin{itemize}[leftmargin=2.6em, labelsep=0.4em, itemsep=2pt, topsep=3pt]
\item[\optlabel{A}] The density matrix of 7 qubits needs exactly 4.10 KB
\item[\optlabel{B}] Each qubit needs $16 \times 7 = 112$ B and there are 36 of them
\item[\optlabel{C}] The state needs $2048$ B and the notebook reserves one extra copy (for reading the full state), giving $4096$ B
\item[\optlabel{D}] Custom gates add $2048$ B of matrix storage per qubit
\end{itemize}
\end{qitem}
\qsep

\needspace{7\baselineskip}
\begin{qitem}{34}
Preparing an arbitrary normalised $n$-qubit amplitude vector $\ket b$ from $\ket{0\dots0}$ with one- and two-qubit gates requires in general
\par\smallskip\noindent\begin{tabularx}{\linewidth}{@{}*{2}{L}@{}}
\optlabel{A}\,$O(n)$ gates & \optlabel{B}\,$O(\log n)$ gates \\[3pt]
\optlabel{C}\,a single gate per qubit & \optlabel{D}\,$O(2^n)$ gates
\end{tabularx}
\end{qitem}
\qsep

\needspace{7\baselineskip}
\begin{qitem}{35}
\textbf{Assertion (A):} In HHL, measuring the clock register immediately after phase estimation would destroy the superposition over eigencomponents that $\ket x$ needs.\\[1mm]
\textbf{Reason (R):} The inverse QFT is a unitary operation.
\begin{itemize}[leftmargin=2.6em, labelsep=0.4em, itemsep=2pt, topsep=3pt]
\item[\optlabel{A}] Both Assertion (A) and Reason (R) are true, and R is the correct explanation of A
\item[\optlabel{B}] Both A and R are true, but R is \emph{not} the correct explanation of A
\item[\optlabel{C}] A is true, but R is false
\item[\optlabel{D}] A is false, but R is true
\end{itemize}
\end{qitem}
\qsep

\needspace{7\baselineskip}
\begin{qitem}{36}
A swap test between states $\ket\psi$ and $\ket\phi$ gives the ancilla outcome 0 with probability $P(0) = \tfrac12(1 + |\braket{\psi|\phi}|^2)$. If $P(0) = 0.9$ is measured, the overlap
$|\braket{\psi|\phi}|^2$ is
\par\smallskip\noindent\begin{tabularx}{\linewidth}{@{}*{4}{L}@{}}
\optlabel{A}\,$0.8$ & \optlabel{B}\,$0.95$ & \optlabel{C}\,$0.9$ & \optlabel{D}\,$0.4$
\end{tabularx}
\end{qitem}
\qsep

\needspace{20\baselineskip}
\begin{passagebox}{Paragraph for Questions 37 and 38}
Quantum phase estimation with $t = 2$ counting qubits is applied to an eigenstate $\ket u$ of a unitary $U$, $U\ket u = e^{2\pi i\varphi}\ket u$. The counting
register ends in $\sum_y\alpha_y\ket y$ with $\alpha_y = \frac14\sum_{j=0}^{3}e^{2\pi i j(\varphi - y/4)}$, and the outcome $y$ estimates $\varphi \approx y/4$.
\end{passagebox}
\needspace{7\baselineskip}
\begin{qitem}{37}
For $\varphi = \tfrac14$, the measured counting register is
\par\smallskip\noindent\begin{tabularx}{\linewidth}{@{}*{2}{L}@{}}
\optlabel{A}\,$\ket{10}$ with probability 1 & \optlabel{B}\,$\ket{01}$ with probability 1 \\[3pt]
\optlabel{C}\,$\ket{01}$ with probability $\tfrac14$ & \optlabel{D}\,$\ket{00}$ with probability 1
\end{tabularx}
\end{qitem}
\qsep

\needspace{7\baselineskip}
\begin{qitem}{38}
For $\varphi = \tfrac13$, the most likely outcome and its probability are
\par\smallskip\noindent\begin{tabularx}{\linewidth}{@{}*{2}{L}@{}}
\optlabel{A}\,all four outcomes, $0.25$ each & \optlabel{B}\,$\ket{01}$, exactly $1$ \\[3pt]
\optlabel{C}\,$\ket{10}$, about $0.70$ & \optlabel{D}\,$\ket{01}$, about $0.70$
\end{tabularx}
\end{qitem}
\qsep

\needspace{7\baselineskip}
\begin{qitem}{39}
\textbf{Statement I:} The inverse QFT is the adjoint of the QFT, so applying it after the QFT returns the original state.\\[1mm]
\textbf{Statement II:} The QFT on $n$ qubits requires a number of gates exponential in $n$.
\begin{itemize}[leftmargin=2.6em, labelsep=0.4em, itemsep=2pt, topsep=3pt]
\item[\optlabel{A}] Both Statement I and Statement II are true
\item[\optlabel{B}] Both Statement I and Statement II are false
\item[\optlabel{C}] Statement I is true, but Statement II is false
\item[\optlabel{D}] Statement I is false, but Statement II is true
\end{itemize}
\end{qitem}
\qsep

\needspace{7\baselineskip}
\begin{qitem}{40}
With $n_c = 4$ and $t = \pi/8$, the matrix $A$ is replaced by $A + I$. How do the ideal clock values change?
\begin{itemize}[leftmargin=2.6em, labelsep=0.4em, itemsep=2pt, topsep=3pt]
\item[\optlabel{A}] each clock value increases by $1$
\item[\optlabel{B}] each clock value increases by $\pi/8$
\item[\optlabel{C}] they are unchanged, because $I$ only adds a global phase
\item[\optlabel{D}] each clock value doubles
\end{itemize}
\end{qitem}
\qsep

\newpage

\begin{sectionbanner}SECTION C \enspace$\cdot$\enspace HHL Algorithm Fundamentals \hfill Q.41--60\end{sectionbanner}

{\small \textbf{Syllabus:} registers, state preparation, phase estimation with $e^{iAt}$, eigenvalue inversion and the constant $C$, controlled rotations, uncomputation, post-selection, norm recovery and the notebook's circuit and code.\\ \textbf{Formats:} Single-correct MCQ (Q.41--57); Matching lists (Q.46--58); Assertion--Reason (Q.49); Statement I / II (Q.52); Paragraph based (Q.59--60). Each question: $+3$ / $0$ / $-1$.}
\par\medskip

\needspace{7\baselineskip}
\begin{qitem}{41}
After phase estimation and the controlled rotation (before uncomputation), the ideal HHL state is (system $\otimes$ clock $\otimes$ ancilla)
\par\smallskip\noindent\begin{tabularx}{\linewidth}{@{}*{2}{L}@{}}
\optlabel{A}\,$\sum_j\beta_j\ket{u_j}\ket{0}\Big(\sqrt{1 - \tfrac{C^2}{\lambda_j^2}}\,\ket0 + \tfrac{C}{\lambda_j}\ket1\Big)$ & \optlabel{B}\,$\sum_j\tfrac{\beta_j}{\lambda_j}\ket{u_j}\ket0\ket1$ \\[3pt]
\optlabel{C}\,$\sum_j\beta_j\ket{u_j}\ket{k_j}\Big(\sqrt{1 - C^2\lambda_j^2}\,\ket0 + C\lambda_j\ket1\Big)$ & \optlabel{D}\,$\sum_j\beta_j\ket{u_j}\ket{k_j}\Big(\sqrt{1 - \tfrac{C^2}{\lambda_j^2}}\,\ket0 + \tfrac{C}{\lambda_j}\ket1\Big)$
\end{tabularx}
\end{qitem}
\qsep

\needspace{7\baselineskip}
\begin{qitem}{42}
Using the notebook's layout (clock + system + one ancilla), an $8\times8$ Hermitian system with $n_c = 5$ clock qubits needs how many qubits?
\par\smallskip\noindent\begin{tabularx}{\linewidth}{@{}*{4}{L}@{}}
\optlabel{A}\,$14$ & \optlabel{B}\,$9$ & \optlabel{C}\,$13$ & \optlabel{D}\,$8$
\end{tabularx}
\end{qitem}
\qsep

\needspace{7\baselineskip}
\begin{qitem}{43}
The notebook sets $C = 2\pi/(2^{n_c}t)$, the smallest non-zero eigenvalue magnitude the clock can represent. Why must $C \le |\lambda_j|$ for every eigenvalue?
\begin{itemize}[leftmargin=2.6em, labelsep=0.4em, itemsep=2pt, topsep=3pt]
\item[\optlabel{A}] The rotation amplitude $C/\lambda_j$ must satisfy $|C/\lambda_j| \le 1$ so that $\arcsin(C/\lambda_j)$ exists
\item[\optlabel{B}] Otherwise phase estimation needs more counting qubits than system qubits
\item[\optlabel{C}] $C$ must be the largest eigenvalue so that the success probability is 1
\item[\optlabel{D}] Otherwise $e^{iAt}$ would not be unitary
\end{itemize}
\end{qitem}
\qsep

\needspace{7\baselineskip}
\begin{qitem}{44}
With $C = 1$, the eigenvalue-inversion rotation for the clock value of $\lambda = 2$ is $R_Y(\theta)$ with
\par\smallskip\noindent\begin{tabularx}{\linewidth}{@{}*{4}{L}@{}}
\optlabel{A}\,$\theta = \tfrac\pi2$ & \optlabel{B}\,$\theta = \arcsin\tfrac12 = \tfrac\pi6$ & \optlabel{C}\,$\theta = 2\arcsin2$ & \optlabel{D}\,$\theta = 2\arcsin\tfrac12 = \tfrac\pi3$
\end{tabularx}
\end{qitem}
\qsep

\needspace{7\baselineskip}
\begin{qitem}{45}
For the default system ($\beta_j = \tfrac12$ for the eigenvalues $1, 2, 3, 4$; $C = 1$), the ideal success probability $P_{\text{success}} = \sum_j|\beta_j|^2C^2/\lambda_j^2$ is
\par\smallskip\noindent\begin{tabularx}{\linewidth}{@{}*{4}{L}@{}}
\optlabel{A}\,$1.4236$ & \optlabel{B}\,$0.25$ & \optlabel{C}\,$0.3559$ & \optlabel{D}\,$0.5208$
\end{tabularx}
\end{qitem}
\qsep

\needspace{7\baselineskip}
\begin{qitem}{46}
Match each stage of the notebook's HHL circuit (List-I) with what it does to the registers (List-II).
\par\smallskip\noindent{\renewcommand{\arraystretch}{1.35}%
\begin{tabularx}{\linewidth}{@{}c>{\raggedright\arraybackslash}p{0.33\linewidth}@{\hspace{1.2em}}cL@{}}
\toprule
\multicolumn{2}{@{}l}{\textbf{List-I}} & \multicolumn{2}{l@{}}{\textbf{List-II}} \\ \midrule
\textbf{(P)} & amplitude initialisation & \textbf{(1)} & $\ket{u_j}\ket{k_j}\ket0 \to \ket{u_j}\ket{k_j}\big(\sqrt{1 - C^2/\lambda_j^2}\ket0 + \tfrac{C}{\lambda_j}\ket1\big)$ \\
\textbf{(Q)} & phase estimation & \textbf{(2)} & returns the clock register to $\ket{0\dots0}$ so that it is no longer entangled with the system \\
\textbf{(R)} & eigenvalue inversion & \textbf{(3)} & prepares $\ket{0\dots0}_{\text{clock}}\ket b\ket0_{\text{anc}}$ \\
\textbf{(S)} & uncomputation & \textbf{(4)} & $\ket b\ket0 \to \sum_j\beta_j\ket{u_j}\ket{k_j}$ \\
\bottomrule
\end{tabularx}}
\par\smallskip\noindent\begin{tabularx}{\linewidth}{@{}*{2}{L}@{}}
\optlabel{A}\,P$\to$3, Q$\to$4, R$\to$1, S$\to$2 & \optlabel{B}\,P$\to$3, Q$\to$4, R$\to$2, S$\to$1 \\[3pt]
\optlabel{C}\,P$\to$4, Q$\to$3, R$\to$1, S$\to$2 & \optlabel{D}\,P$\to$3, Q$\to$1, R$\to$4, S$\to$2
\end{tabularx}
\end{qitem}
\qsep

\needspace{7\baselineskip}
\begin{qitem}{47}
With \texttt{EVOLUTION\_TIME = "auto"} the notebook chooses $t = 2\pi/(2^{n_c}|\lambda|_{\min})$ whenever $\kappa \le 2^{n_c - 1} - 1$. For $n_c = 4$ and eigenvalues $0.5, 1, 2, 3$, it picks
\par\smallskip\noindent\begin{tabularx}{\linewidth}{@{}*{2}{L}@{}}
\optlabel{A}\,$t = \pi/8$ and $C = 1$ & \optlabel{B}\,$t = 2\pi/3$ and $C = 3$ \\[3pt]
\optlabel{C}\,$t = \pi/4$ and $C = 1$ & \optlabel{D}\,$t = \pi/4$ and $C = 0.5$
\end{tabularx}
\end{qitem}
\qsep

\needspace{7\baselineskip}
\begin{qitem}{48}
Why does HHL uncompute phase estimation after the eigenvalue inversion instead of simply measuring the ancilla?
\begin{itemize}[leftmargin=2.6em, labelsep=0.4em, itemsep=2pt, topsep=3pt]
\item[\optlabel{A}] Uncomputation doubles the success probability
\item[\optlabel{B}] Uncomputation is needed to make $e^{iAt}$ unitary
\item[\optlabel{C}] It converts the amplitudes $C/\lambda_j$ into probabilities
\item[\optlabel{D}] Without it the clock stays entangled with the system, so the post-selected system register would be a mixture instead of $\ket x$ (unless the clock is also post-selected on a fixed value)
\end{itemize}
\end{qitem}
\qsep

\needspace{7\baselineskip}
\begin{qitem}{49}
\textbf{Assertion (A):} From the simulated state vector, the notebook recovers the full solution $x$ including its norm.\\[1mm]
\textbf{Reason (R):} In the post-selected branch (clock $= 0$, ancilla $= 1$) the system amplitudes are $\tfrac{C}{\|b\|}\,x$, so $x = \tfrac{\|b\|}{C}\times(\text{amplitudes})$.
\begin{itemize}[leftmargin=2.6em, labelsep=0.4em, itemsep=2pt, topsep=3pt]
\item[\optlabel{A}] Both Assertion (A) and Reason (R) are true, and R is the correct explanation of A
\item[\optlabel{B}] Both A and R are true, but R is \emph{not} the correct explanation of A
\item[\optlabel{C}] A is true, but R is false
\item[\optlabel{D}] A is false, but R is true
\end{itemize}
\end{qitem}
\qsep

\needspace{7\baselineskip}
\begin{qitem}{50}
The notebook's eigenvalue inversion applies one multi-controlled rotation for every non-zero clock value. With $n_c$ clock qubits this is
\par\smallskip\noindent\begin{tabularx}{\linewidth}{@{}*{2}{L}@{}}
\optlabel{A}\,$n_c$ rotations & \optlabel{B}\,$2^{n_c} - 1$ rotations ($15$ for $n_c = 4$) \\[3pt]
\optlabel{C}\,one rotation per eigenvalue of $A$ & \optlabel{D}\,$2^{n_c}$ rotations, including $k = 0$
\end{tabularx}
\end{qitem}
\qsep

\needspace{7\baselineskip}
\begin{qitem}{51}
With $n_c = 4$ and $t = \pi/8$, the notebook assigns to the clock value $k = 12$ the eigenvalue
\par\smallskip\noindent\begin{tabularx}{\linewidth}{@{}*{4}{L}@{}}
\optlabel{A}\,$12$ & \optlabel{B}\,$-4$ & \optlabel{C}\,$4$ & \optlabel{D}\,$-12$
\end{tabularx}
\end{qitem}
\qsep

\needspace{7\baselineskip}
\begin{qitem}{52}
\textbf{Statement I:} If some $\lambda_jt\,2^{n_c}/(2\pi)$ is not an integer, phase estimation spreads that component over several clock values and the fidelity of $\ket x$ drops below 1.\\[1mm]
\textbf{Statement II:} The notebook's \texttt{"auto"} evolution time guarantees exact clock values for every matrix.
\begin{itemize}[leftmargin=2.6em, labelsep=0.4em, itemsep=2pt, topsep=3pt]
\item[\optlabel{A}] Both Statement I and Statement II are true
\item[\optlabel{B}] Both Statement I and Statement II are false
\item[\optlabel{C}] Statement I is true, but Statement II is false
\item[\optlabel{D}] Statement I is false, but Statement II is true
\end{itemize}
\end{qitem}
\qsep

\needspace{7\baselineskip}
\begin{qitem}{53}
For the notebook's $2\times2$ example ($\lambda = \tfrac23, \tfrac43$) with $n_c = 4$ and \texttt{"auto"} time, the values of $t$, $C$ and the clock values are
\par\smallskip\noindent\begin{tabularx}{\linewidth}{@{}*{2}{L}@{}}
\optlabel{A}\,$t = \tfrac{3\pi}{16}$, $C = \tfrac43$, clock values $1$ and $2$ & \optlabel{B}\,$t = \tfrac{3\pi}{8}$, $C = \tfrac23$, clock values $2$ and $4$ \\[3pt]
\optlabel{C}\,$t = \tfrac\pi8$, $C = 1$, clock values $\tfrac23$ and $\tfrac43$ & \optlabel{D}\,$t = \tfrac{3\pi}{16}$, $C = \tfrac23$, clock values $1$ and $2$
\end{tabularx}
\end{qitem}
\qsep

\needspace{7\baselineskip}
\begin{qitem}{54}
In the sampling read-out the notebook keeps a shot only if
\begin{itemize}[leftmargin=2.6em, labelsep=0.4em, itemsep=2pt, topsep=3pt]
\item[\optlabel{A}] the clock bits encode an eigenvalue
\item[\optlabel{B}] the ancilla bit is 0
\item[\optlabel{C}] all clock bits are 0 and the ancilla bit is 1
\item[\optlabel{D}] the system bits are all 0
\end{itemize}
\end{qitem}
\qsep

\needspace{7\baselineskip}
\begin{qitem}{55}
The ideal success probability is $p = 0.3559$ and the notebook uses $20\,000$ shots. The standard error of the measured fraction of accepted shots is about
\par\smallskip\noindent\begin{tabularx}{\linewidth}{@{}*{4}{L}@{}}
\optlabel{A}\,$0.0034$ & \optlabel{B}\,$0.0001$ & \optlabel{C}\,$0.0119$ & \optlabel{D}\,$0.048$
\end{tabularx}
\end{qitem}
\qsep

\needspace{7\baselineskip}
\begin{qitem}{56}
On hardware, the post-selected system bitstrings of HHL directly estimate
\begin{itemize}[leftmargin=2.6em, labelsep=0.4em, itemsep=2pt, topsep=3pt]
\item[\optlabel{A}] $\|x\|$ only
\item[\optlabel{B}] $|x_i|^2/\|x\|^2$ only --- signs/phases and the norm need further measurements
\item[\optlabel{C}] $x_i/\|x\|$ including its sign
\item[\optlabel{D}] every component $x_i$ including its sign
\end{itemize}
\end{qitem}
\qsep

\needspace{7\baselineskip}
\begin{qitem}{57}
The recorded circuit statistics of the default HHL circuit are\\
\hspace*{1em}\texttt{h: 16, ccustom\_unitary\_constant: 8, swap: 4, cr1: 12, x: 56, ccccry: 15}.\\
How many multi-qubit gates is that?
\par\smallskip\noindent\begin{tabularx}{\linewidth}{@{}*{4}{L}@{}}
\optlabel{A}\,$111$ & \optlabel{B}\,$39$ & \optlabel{C}\,$23$ & \optlabel{D}\,$95$
\end{tabularx}
\end{qitem}
\qsep

\needspace{7\baselineskip}
\begin{qitem}{58}
Match each HHL ingredient (List-I) with the code that implements it in \texttt{HHL.ipynb} (List-II).
\par\smallskip\noindent{\renewcommand{\arraystretch}{1.35}%
\begin{tabularx}{\linewidth}{@{}c>{\raggedright\arraybackslash}p{0.33\linewidth}@{\hspace{1.2em}}cL@{}}
\toprule
\multicolumn{2}{@{}l}{\textbf{List-I}} & \multicolumn{2}{l@{}}{\textbf{List-II}} \\ \midrule
\textbf{(P)} & preparation of $\ket0\ket b\ket0$ & \textbf{(1)} & \texttt{kernel.cry(theta, clock, ancilla)} between $X$ gates \\
\textbf{(Q)} & controlled powers $U^{2^j}$ & \textbf{(2)} & \texttt{postselected\_amplitudes(state\_vector)} \\
\textbf{(R)} & eigenvalue inversion & \textbf{(3)} & \texttt{kernel.qalloc(bit\_reverse\_order(initial, n\_tot))} \\
\textbf{(S)} & extraction of the post-selected amplitudes & \textbf{(4)} & \texttt{cudaq.register\_operation(name, matrix\_power(U, 2**j))} \\
\bottomrule
\end{tabularx}}
\par\smallskip\noindent\begin{tabularx}{\linewidth}{@{}*{2}{L}@{}}
\optlabel{A}\,P$\to$4, Q$\to$3, R$\to$1, S$\to$2 & \optlabel{B}\,P$\to$3, Q$\to$1, R$\to$4, S$\to$2 \\[3pt]
\optlabel{C}\,P$\to$3, Q$\to$4, R$\to$1, S$\to$2 & \optlabel{D}\,P$\to$3, Q$\to$4, R$\to$2, S$\to$1
\end{tabularx}
\end{qitem}
\qsep

\needspace{20\baselineskip}
\begin{passagebox}{Paragraph for Questions 59 and 60}
HHL in the notebook's convention is applied to $A = \operatorname{diag}(1, 2)$ and $b = \tfrac{1}{\sqrt2}(1, 1)$ with $n_c = 3$ clock qubits and $t = \pi/4$, so
$U = e^{iAt}$, the clock value of an eigenvalue $\lambda$ is $k = \lambda t\,2^{n_c}/(2\pi) = \lambda$, and $C = 2\pi/(2^{n_c}t) = 1$. The eigenvalue inversion applies
$R_Y(2\arcsin(C/\lambda_k))$ to the ancilla for each clock value; the system register has one qubit.
\end{passagebox}
\needspace{7\baselineskip}
\begin{qitem}{59}
After phase estimation the state of system $\otimes$ clock (ancilla still $\ket0$) is
\par\smallskip\noindent\begin{tabularx}{\linewidth}{@{}*{2}{L}@{}}
\optlabel{A}\,$\tfrac{1}{\sqrt2}\big(\ket0\ket{001} + \ket1\ket{011}\big)$ & \optlabel{B}\,$\tfrac{1}{\sqrt2}\big(\ket0\ket{100} + \ket1\ket{010}\big)$ \\[3pt]
\optlabel{C}\,$\tfrac{1}{\sqrt2}\big(\ket0\ket{001} + \ket1\ket{010}\big)$ & \optlabel{D}\,$\tfrac{1}{\sqrt2}\big(\ket0 + \ket1\big)\ket{001}$
\end{tabularx}
\end{qitem}
\qsep

\needspace{7\baselineskip}
\begin{qitem}{60}
The success probability and the normalised output state $\ket x$ are
\par\smallskip\noindent\begin{tabularx}{\linewidth}{@{}*{2}{L}@{}}
\optlabel{A}\,$0.625$ and $\tfrac{1}{\sqrt5}(2, 1)$ & \optlabel{B}\,$0.625$ and $\tfrac{1}{\sqrt2}(1, 1)$ \\[3pt]
\optlabel{C}\,$0.75$ and $\tfrac{1}{\sqrt5}(2, 1)$ & \optlabel{D}\,$1$ and $\tfrac{1}{\sqrt5}(1, 2)$
\end{tabularx}
\end{qitem}
\qsep

\newpage

\begin{sectionbanner}SECTION D \enspace$\cdot$\enspace Advanced HHL and Implementation \hfill Q.61--80\end{sectionbanner}

{\small \textbf{Syllabus:} complexity, condition number and success probability, amplitude amplification, read-out cost, clock precision and aliasing, embedding and padding, Hamiltonian simulation, newer algorithms, backends, recorded results and limitations.\\ \textbf{Formats:} Single-correct MCQ (Q.61--78); Matching lists (Q.64--72); Assertion--Reason (Q.69); Statement I / II (Q.74); Paragraph based (Q.79--80). Each question: $+3$ / $0$ / $-1$.}
\par\medskip

\needspace{7\baselineskip}
\begin{qitem}{61}
The original HHL analysis (sparse $d\times d$ matrix with sparsity $s$, condition number $\kappa$, precision $\varepsilon$) gives a running time of
\par\smallskip\noindent\begin{tabularx}{\linewidth}{@{}*{2}{L}@{}}
\optlabel{A}\,$O(\log d)$, independent of $\kappa$ and $\varepsilon$ & \optlabel{B}\,$O(d\,s\sqrt\kappa\log(1/\varepsilon))$ \\[3pt]
\optlabel{C}\,$O(d^3)$ & \optlabel{D}\,$\tilde O\!\left(\log(d)\,s^2\kappa^2/\varepsilon\right)$
\end{tabularx}
\end{qitem}
\qsep

\needspace{7\baselineskip}
\begin{qitem}{62}
HHL succeeds (ancilla $= 1$) with probability $p = 0.01$. On average, how many repetitions are needed with plain repetition, and with amplitude amplification?
\par\smallskip\noindent\begin{tabularx}{\linewidth}{@{}*{2}{L}@{}}
\optlabel{A}\,about $100$ and about $10$ & \optlabel{B}\,about $100$ and about $1$ \\[3pt]
\optlabel{C}\,about $1$ and about $10$ & \optlabel{D}\,about $10$ and about $100$
\end{tabularx}
\end{qitem}
\qsep

\needspace{7\baselineskip}
\begin{qitem}{63}
To estimate all $d = 1024$ probabilities $|x_i|^2$ of $\ket x$ to absolute precision $\varepsilon = 0.01$ by sampling, roughly how many post-selected shots are needed?
\par\smallskip\noindent\begin{tabularx}{\linewidth}{@{}*{4}{L}@{}}
\optlabel{A}\,about $10^{14}$ & \optlabel{B}\,about $10$ & \optlabel{C}\,about $10^7$ & \optlabel{D}\,about $10^3$
\end{tabularx}
\end{qitem}
\qsep

\needspace{7\baselineskip}
\begin{qitem}{64}
For exact clock values the notebook needs $\kappa \le 2^{n_c - 1} - 1$. Match each condition number (List-I) with the smallest suitable number of clock qubits $n_c$
(List-II).
\par\smallskip\noindent{\renewcommand{\arraystretch}{1.35}%
\begin{tabularx}{\linewidth}{@{}c>{\raggedright\arraybackslash}p{0.33\linewidth}@{\hspace{1.2em}}cL@{}}
\toprule
\multicolumn{2}{@{}l}{\textbf{List-I}} & \multicolumn{2}{l@{}}{\textbf{List-II}} \\ \midrule
\textbf{(P)} & $\kappa = 3$ & \textbf{(1)} & $6$ \\
\textbf{(Q)} & $\kappa = 7$ & \textbf{(2)} & $3$ \\
\textbf{(R)} & $\kappa = 20$ & \textbf{(3)} & $7$ \\
\textbf{(S)} & $\kappa = 40$ & \textbf{(4)} & $4$ \\
\bottomrule
\end{tabularx}}
\par\smallskip\noindent\begin{tabularx}{\linewidth}{@{}*{2}{L}@{}}
\optlabel{A}\,P$\to$2, Q$\to$4, R$\to$3, S$\to$1 & \optlabel{B}\,P$\to$2, Q$\to$1, R$\to$4, S$\to$3 \\[3pt]
\optlabel{C}\,P$\to$4, Q$\to$2, R$\to$1, S$\to$3 & \optlabel{D}\,P$\to$2, Q$\to$4, R$\to$1, S$\to$3
\end{tabularx}
\end{qitem}
\qsep

\needspace{7\baselineskip}
\begin{qitem}{65}
With \texttt{NUM\_CLOCK\_QUBITS = 8}, how many multi-controlled eigenvalue-inversion rotations does the notebook's circuit contain, and how many qubits does the
default $4\times4$ system then use?
\par\smallskip\noindent\begin{tabularx}{\linewidth}{@{}*{2}{L}@{}}
\optlabel{A}\,$255$ rotations, $11$ qubits & \optlabel{B}\,$255$ rotations, $14$ qubits \\[3pt]
\optlabel{C}\,$256$ rotations, $12$ qubits & \optlabel{D}\,$8$ rotations, $11$ qubits
\end{tabularx}
\end{qitem}
\qsep

\needspace{7\baselineskip}
\begin{qitem}{66}
With \texttt{MATRIX\_A = [[1, 2], [0, 1]]} (not Hermitian) and $n_c = 4$, the notebook's circuit uses how many qubits, and where is $x$ in the solution of the embedded system?
\begin{itemize}[leftmargin=2.6em, labelsep=0.4em, itemsep=2pt, topsep=3pt]
\item[\optlabel{A}] $7$ qubits; $x$ is the first half of the solution
\item[\optlabel{B}] $7$ qubits; $x$ is the second half of the solution
\item[\optlabel{C}] $8$ qubits; $x$ is the second half
\item[\optlabel{D}] $6$ qubits; $x$ is the whole solution
\end{itemize}
\end{qitem}
\qsep

\needspace{7\baselineskip}
\begin{qitem}{67}
A $3\times3$ Hermitian matrix with eigenvalues $5, 10, 20$ is padded by the notebook to $4\times4$ with an identity block. What happens to $\kappa$?
\begin{itemize}[leftmargin=2.6em, labelsep=0.4em, itemsep=2pt, topsep=3pt]
\item[\optlabel{A}] it grows from $4$ to $20$, because the padding adds the eigenvalue $1$
\item[\optlabel{B}] it stays $4$
\item[\optlabel{C}] it becomes $5$
\item[\optlabel{D}] it becomes $1$
\end{itemize}
\end{qitem}
\qsep

\needspace{7\baselineskip}
\begin{qitem}{68}
``Dequantised'' algorithms (Tang, 2018 and later) showed that for some quantum linear-algebra tasks
\begin{itemize}[leftmargin=2.6em, labelsep=0.4em, itemsep=2pt, topsep=3pt]
\item[\optlabel{A}] HHL can run without phase estimation
\item[\optlabel{B}] the condition number no longer affects quantum algorithms
\item[\optlabel{C}] classical computers can prepare $\ket x$ on a quantum computer
\item[\optlabel{D}] classical randomized algorithms with comparable sampling access achieve similar poly-logarithmic dependence on $d$ for low-rank matrices
\end{itemize}
\end{qitem}
\qsep

\needspace{7\baselineskip}
\begin{qitem}{69}
\textbf{Assertion (A):} HHL gives an exponential speed-up over classical algorithms for every linear system $Ax = b$.\\[1mm]
\textbf{Reason (R):} For $s$-sparse, well-conditioned $A$ with efficiently preparable $\ket b$, HHL prepares $\ket x$ in time polynomial in $\log d$, $s$, $\kappa$ and $1/\varepsilon$.
\begin{itemize}[leftmargin=2.6em, labelsep=0.4em, itemsep=2pt, topsep=3pt]
\item[\optlabel{A}] Both Assertion (A) and Reason (R) are true, and R is the correct explanation of A
\item[\optlabel{B}] Both A and R are true, but R is \emph{not} the correct explanation of A
\item[\optlabel{C}] A is true, but R is false
\item[\optlabel{D}] A is false, but R is true
\end{itemize}
\end{qitem}
\qsep

\needspace{7\baselineskip}
\begin{qitem}{70}
Take $C = \lambda_{\min}$ and a matrix with $\kappa = 20$. If $\ket b$ is the eigenvector of $\lambda_{\max}$, the success probability is
\par\smallskip\noindent\begin{tabularx}{\linewidth}{@{}*{4}{L}@{}}
\optlabel{A}\,$0.5$ & \optlabel{B}\,$1/\kappa^2 = 0.0025$ & \optlabel{C}\,$1$ & \optlabel{D}\,$1/\kappa = 0.05$
\end{tabularx}
\end{qitem}
\qsep

\needspace{7\baselineskip}
\begin{qitem}{71}
With $n_c = 4$ and $t = \pi/8$ (clock value $k = \lambda$), how is an eigenvalue $\lambda = 9$ treated?
\begin{itemize}[leftmargin=2.6em, labelsep=0.4em, itemsep=2pt, topsep=3pt]
\item[\optlabel{A}] Phase estimation fails with an error
\item[\optlabel{B}] It is read as $7$, the largest positive value
\item[\optlabel{C}] It aliases: clock value $9$ is read as the signed value $-7$, so the wrong (negative) $1/\lambda$ is applied
\item[\optlabel{D}] It is read correctly as $9$
\end{itemize}
\end{qitem}
\qsep

\needspace{7\baselineskip}
\begin{qitem}{72}
Match each quantum linear-system algorithm (List-I) with its key idea (List-II).
\par\smallskip\noindent{\renewcommand{\arraystretch}{1.35}%
\begin{tabularx}{\linewidth}{@{}c>{\raggedright\arraybackslash}p{0.33\linewidth}@{\hspace{1.2em}}cL@{}}
\toprule
\multicolumn{2}{@{}l}{\textbf{List-I}} & \multicolumn{2}{l@{}}{\textbf{List-II}} \\ \midrule
\textbf{(P)} & HHL (2009) & \textbf{(1)} & implement a polynomial approximation of $1/x$ on the singular values of a block-encoded $A$ \\
\textbf{(Q)} & Ambainis (2012), variable-time amplitude amplification & \textbf{(2)} & phase estimation + controlled rotation $C/\lambda$ + post-selection; cost $\kappa^2/\varepsilon$ \\
\textbf{(R)} & Childs--Kothari--Somma (2017) & \textbf{(3)} & amplify different eigenvalue branches at different times: $\kappa$ dependence close to linear ($\kappa\log^3\kappa$) \\
\textbf{(S)} & QSVT-based solvers & \textbf{(4)} & linear combination of unitaries (Fourier/Chebyshev series of $1/x$): $\mathrm{poly}\log(1/\varepsilon)$ precision dependence \\
\bottomrule
\end{tabularx}}
\par\smallskip\noindent\begin{tabularx}{\linewidth}{@{}*{2}{L}@{}}
\optlabel{A}\,P$\to$2, Q$\to$4, R$\to$3, S$\to$1 & \optlabel{B}\,P$\to$3, Q$\to$2, R$\to$4, S$\to$1 \\[3pt]
\optlabel{C}\,P$\to$2, Q$\to$3, R$\to$4, S$\to$1 & \optlabel{D}\,P$\to$2, Q$\to$3, R$\to$1, S$\to$4
\end{tabularx}
\end{qitem}
\qsep

\needspace{7\baselineskip}
\begin{qitem}{73}
With \texttt{SIMULATION\_METHOD = "density\_matrix"} the notebook multiplies its recovered $x$ by a phase factor before comparing with NumPy. Why?
\begin{itemize}[leftmargin=2.6em, labelsep=0.4em, itemsep=2pt, topsep=3pt]
\item[\optlabel{A}] A density matrix $\ket\psi\bra\psi$ carries no global phase, so the vector extracted from it has an arbitrary phase that must be aligned
\item[\optlabel{B}] The density-matrix simulator uses fp32 precision
\item[\optlabel{C}] The density-matrix simulator adds noise that rotates the phase
\item[\optlabel{D}] Post-selection on a density matrix flips the sign of every component
\end{itemize}
\end{qitem}
\qsep

\needspace{7\baselineskip}
\begin{qitem}{74}
\textbf{Statement I:} \texttt{HHL.ipynb} refuses hardware backends, because exact matrix exponentials as custom gates and amplitude initialisation cannot run on hardware targets.\\[1mm]
\textbf{Statement II:} The controlled powers $U^{2^j}$ are registered as custom unitary matrices with \texttt{cudaq.register\_operation}.
\begin{itemize}[leftmargin=2.6em, labelsep=0.4em, itemsep=2pt, topsep=3pt]
\item[\optlabel{A}] Both Statement I and Statement II are true
\item[\optlabel{B}] Both Statement I and Statement II are false
\item[\optlabel{C}] Statement I is true, but Statement II is false
\item[\optlabel{D}] Statement I is false, but Statement II is true
\end{itemize}
\end{qitem}
\qsep

\needspace{7\baselineskip}
\begin{qitem}{75}
Run with only \texttt{NUM\_CLOCK\_QUBITS = 3} on the default system ($\kappa = 4 > 3$ levels), the notebook rescales $t$ so that $\lambda_{\max} = 4$ maps to clock value 3. The ideal
clock values of $\lambda = 1, 2, 3, 4$ and the effect are
\begin{itemize}[leftmargin=2.6em, labelsep=0.4em, itemsep=2pt, topsep=3pt]
\item[\optlabel{A}] $1, 2, 3, 4$ --- still exact
\item[\optlabel{B}] $0.75, 1.5, 2.25, 3$ --- non-integer values are rounded by phase estimation, so the fidelity drops below 1
\item[\optlabel{C}] $3, 6, 9, 12$ --- all eigenvalues alias
\item[\optlabel{D}] $0.75, 1.5, 2.25, 3$ --- the fidelity stays exactly 1
\end{itemize}
\end{qitem}
\qsep

\needspace{7\baselineskip}
\begin{qitem}{76}
With \texttt{NUM\_CLOCK\_QUBITS = 8} the default circuit has 11 qubits. The fp64 state vector needs ... bytes, and the notebook's memory line (state + one host copy) shows about ...
\par\smallskip\noindent\begin{tabularx}{\linewidth}{@{}*{4}{L}@{}}
\optlabel{A}\,$32\,768$; $32.8$ kB & \optlabel{B}\,$32\,768$; $65.5$ kB & \optlabel{C}\,$4\,194\,304$; $8.4$ MB & \optlabel{D}\,$2\,048$; $4.1$ kB
\end{tabularx}
\end{qitem}
\qsep

\needspace{7\baselineskip}
\begin{qitem}{77}
Which statement correctly describes how \texttt{HHL.ipynb} implements $U = e^{iAt}$?
\begin{itemize}[leftmargin=2.6em, labelsep=0.4em, itemsep=2pt, topsep=3pt]
\item[\optlabel{A}] It uses qubitisation with a block encoding of $A$
\item[\optlabel{B}] It uses first-order Trotterisation of $A$'s Pauli decomposition
\item[\optlabel{C}] It computes $e^{iAt}$ exactly with \texttt{scipy.linalg.expm} and applies each power $U^{2^j}$ as a custom controlled gate; this is exact on simulators but is not a scalable Hamiltonian simulation
\item[\optlabel{D}] It approximates $e^{iAt}$ by $I + iAt$
\end{itemize}
\end{qitem}
\qsep

\needspace{7\baselineskip}
\begin{qitem}{78}
Which use of HHL can, in principle, keep an exponential advantage?
\begin{itemize}[leftmargin=2.6em, labelsep=0.4em, itemsep=2pt, topsep=3pt]
\item[\optlabel{A}] replacing Gaussian elimination for $4\times4$ systems
\item[\optlabel{B}] printing every component of $x$ for a dense $10^6\times10^6$ matrix
\item[\optlabel{C}] solving a dense, ill-conditioned system with $\kappa \sim d$
\item[\optlabel{D}] estimating a quantity such as $x^\dagger Mx$ for a sparse, well-conditioned $A$ with efficiently preparable $\ket b$
\end{itemize}
\end{qitem}
\qsep

\needspace{20\baselineskip}
\begin{passagebox}{Paragraph for Questions 79 and 80}
\emph{Previously recorded default run of \texttt{HHL.ipynb}} (CUDA-Q 0.16.0, target \texttt{qpp-cpu}): $4\times4$ system, $\kappa = 4$, 7 qubits, $\|b\| = 1$, $C = 1$.
Success probability: $0.355903$ from the state, $0.3543$ from $20\,000$ shots ($7086$ accepted). Exact $|x_i|^2/\|x\|^2 = (0.7622, 0.0598, 0.1476, 0.0305)$; from the accepted shots
$(0.7691, 0.0587, 0.1420, 0.0302)$. Fidelity of the state read-out $1.0000000000$, relative error $1.655\times10^{-15}$.
\end{passagebox}
\needspace{7\baselineskip}
\begin{qitem}{79}
The sampled $|x_0|^2/\|x\|^2 = 0.7691$ differs from the exact $0.7622$ by $0.0069$. With $7086$ accepted shots, its binomial standard error is about ..., so the deviation is ...
\begin{itemize}[leftmargin=2.6em, labelsep=0.4em, itemsep=2pt, topsep=3pt]
\item[\optlabel{A}] $0.0051$; about $14$ standard errors --- a bug
\item[\optlabel{B}] $0.0051$; about $1.4$ standard errors --- consistent with shot noise
\item[\optlabel{C}] $0.0034$; about $2$ standard errors
\item[\optlabel{D}] $0.0119$; about $0.6$ standard errors
\end{itemize}
\end{qitem}
\qsep

\needspace{7\baselineskip}
\begin{qitem}{80}
Using $P_{\text{success}} = C^2\|x\|^2/\|b\|^2$, the norm $\|x\|$ estimated from the shot-based success fraction is about
\par\smallskip\noindent\begin{tabularx}{\linewidth}{@{}*{4}{L}@{}}
\optlabel{A}\,$1.68$ & \optlabel{B}\,$0.354$ & \optlabel{C}\,$0.595$ & \optlabel{D}\,$0.125$
\end{tabularx}
\end{qitem}
\qsep

\newpage

\end{document}